\section{Conventions, and the facts we use}\label{sec:conventions}

\subsection{The Weil form}

Inner products are conjugate-linear in the first argument. For $L>0$ let
$I_L=(-L/2,L/2)$, let $P_L$ be multiplication by $\mathbf 1_{I_L}$, and for
$f\in L^2(I_L)$ let $F$ be its extension by zero. Write
\begin{equation}\label{eq:fourier}
 \widehat F(\tau)=\int_\R e^{-i\tau x}F(x)\dd x,\qquad
 U_dF(x)=F(x-d),\qquad
 \norm F^2=\frac1{2\pi}\int_\R|\widehat F(\tau)|^2\dd\tau ,
\end{equation}
and, for $r>0$,
\begin{equation}\label{eq:ngamma}
 \gammak(r)=\frac{e^{-r/2}}{1-e^{-2r}}=\frac{e^{r/2}}{2\sinh r}.
\end{equation}
The localized Weil form is, in the normalization of \cite{WilsonLines},
\begin{multline}\label{eq:target}
 Q_L[f]=\int_0^\infty\big[2\norm F^2-2\re\ip F{U_rF}\big]\gammak(r)\dd r
 +w_0\norm f^2+P_L[f]\\
 -2\!\!\sum_{m\log p<L}\!\!(\log p)\,p^{-m/2}\,\re\ip F{U_{m\log p}F},
\end{multline}
with $w_0=\psi(\frac14)-\log\pi$, $\psi=\Gamma'/\Gamma$, and the pole form
$P_L[f]=2|\int f\cosh\frac x2|^2-2|\int f\sinh\frac x2|^2$. Its spectral form is
\begin{equation}\label{eq:spectral}
 Q_L[f]=\sum_\rho\widehat F(\gamma_\rho)\overline{\widehat F(\overline{\gamma_\rho})},
 \qquad\gamma_\rho=-i(\rho-\tfrac12),
\end{equation}
the sum over the nontrivial zeros of $\zeta$ in symmetric order, and Weil's
criterion, used as input, is $\mathrm{RH}\iff Q_L\geq0$ for every $L$. We write
$m_L$ for the smallest eigenvalue of the (unbounded below-bounded) operator
representing $Q_L$ on $L^2(I_L)$, the \emph{margin}; under RH it is positive
\cite[Section~8]{WilsonLines}, and its recorded values at
$L=\log3,\log5,\log7$ are $5.5\times10^{-8}$, $9.3\times10^{-18}$,
$6.8\times10^{-28}$ \cite{CriticalPath,FrameBound}.

\subsection{The shifted family}

Let
\begin{equation}\label{eq:xi-lambda}
 \Lambda(u)=\pi^{-u/2}\Gamma(u/2)\zeta(u),\qquad
 \xi(u)=\tfrac12u(u-1)\Lambda(u),\qquad s=\tfrac12+p,
\end{equation}
so that $\xi$ is entire, $\xi(u)=\xi(1-u)$, $\xi$ is real on the real axis, and
$\Lambda$ has simple poles at $u=0,1$ with residues $-1,1$ and is otherwise
holomorphic except that it inherits nothing from $\Gamma$ (the trivial zeros of
$\zeta$ cancel the poles of $\Gamma(u/2)$). Throughout,
\begin{equation}\label{eq:ab}
 0<\omega\leq\tfrac12,\qquad a=\tfrac12-\omega,\qquad b=\tfrac12+\omega,
\end{equation}
and the transfer symbol is
\begin{equation}\label{eq:symbol}
 K_\omega(p)=\frac{\xi(s-\omega)}{\xi(s+\omega)}.
\end{equation}
The poles of $K_\omega$ are at the zeros of $\xi(s+\omega)$, $p=\rho-\frac12-\omega$,
all with $\re p\leq\frac12-\omega=a$; its zeros are at $p=\rho-\frac12+\omega$.
The \emph{causal realization} of a symbol $H$ holomorphic on $\re p>\eta_0$ is
the convolution with its inverse Laplace transform on any line $\re p=\eta>\eta_0$,
a distribution $h$ supported on $[0,\infty)$ with $H(p)=\int_0^\infty e^{-pt}h(t)\dd t$;
we write $V_\omega f=k_\omega*f$ for $K_\omega$, and
\begin{equation}\label{eq:compression}
 V_{\omega,L}=P_LV_\omega P_L,\qquad
 D_{\omega,L}=I-V_{\omega,L}^*V_{\omega,L},
\end{equation}
for the compression to the window and its contraction defect. Causal
distributions with a Laplace transform form a commutative convolution algebra
(the algebra $\mathcal D'_+$ of Schwartz), on which the Laplace transform is a
homomorphism; we use this freely.

The \emph{generator} is the causal realization $G_{\omega,L}=P_L(g_\omega*\cdot)P_L$
of the symbol
\begin{equation}\label{eq:generator}
 a_\omega(p)=-\partial_\omega\log K_\omega(p)
 =\frac{\xi'}{\xi}(s+\omega)+\frac{\xi'}{\xi}(s-\omega),
 \qquad
 \frac{\xi'}{\xi}(u)=\frac1u+\frac1{u-1}-\frac12\log\pi+\frac12\psi\Big(\frac u2\Big)+\frac{\zeta'}{\zeta}(u),
\end{equation}
and the \emph{shifted form} is its Hermitian part,
$Q_{\omega,L}[g]=\re\ip g{G_{\omega,L}g}$, with $Q_{0,L}=Q_L$. The flow
identities \cite[Section~2]{WilsonLines} are
\begin{equation}\label{eq:flow}
 \partial_\omega V_{\omega,L}=-G_{\omega,L}V_{\omega,L},\qquad V_{0,L}=I,\qquad
 \frac{d}{d\omega}\norm{V_{\omega,L}f}^2=-2\,Q_{\omega,L}[V_{\omega,L}f],
\end{equation}
so that $\ip f{D_{\omega,L}f}=2\omega Q_L[f]+O(\omega^2)$, which is
\eqref{eq:intro-firstorder}. Two structural facts about $a_\omega$ are used
below. It is \emph{even in $\omega$}, by inspection of \eqref{eq:generator}; at
the kernel level, since a shift of $p$ multiplies a kernel by an exponential,
$g_\omega(t)=2\cosh(\omega t)\,g^+(t)$ with $g^+$ the kernel of
$\frac{\xi'}{\xi}(\frac12+p)$, which is \cite[Theorem~3.1]{WilsonLines}: the
shifted form is the central form with every off-diagonal kernel and every
prime atom multiplied by $\cosh(\omega\cdot)$. And its kernel at $\omega=0$ is
read off term by term from \eqref{eq:generator}:
\begin{equation}\label{eq:g0}
 g_0(t)=2e^{-t/2}+2e^{t/2}-\log\pi\,\delta_0-\fp[2\gammak]-2\sum_{n\geq2}\Lambda(n)\,n^{-1/2}\,\delta_{\log n},
\end{equation}
where the finite part is the causal distribution with Laplace transform
$-\psi(\frac14+\frac p2)$, fixed by the integral representation
$\psi(z)=\int_0^\infty\big(\frac{e^{-t}}t-\frac{e^{-zt}}{1-e^{-t}}\big)dt$:
\begin{equation}\label{eq:fp}
 \big\langle\fp[2\gammak],\phi\big\rangle
 =\lim_{\epsilon\downarrow0}\Big[\int_\epsilon^\infty2\gammak(u)\phi(u)\dd u-\phi(0)\,E_1(2\epsilon)\Big],
 \qquad E_1(x)=\int_x^\infty\frac{e^{-u}}u\dd u .
\end{equation}
(Here $g_0$ is twice $\xi'/\xi$ at $u=\frac12+p$: the constant
$-\frac12\log\pi$ doubles to $-\log\pi\,\delta_0$, and the pole terms
$\frac1u+\frac1{u-1}=\frac1{p+1/2}+\frac1{p-1/2}$ have kernels $e^{-t/2}$ and
$e^{t/2}$, doubled.)

\subsection{Facts about the transfer}

The following are proved in \cite{WilsonLines}; we reprove the two that are
one line, state the rest, and use nothing else from that manuscript.

\begin{proposition}[Unimodularity {\cite[Proposition~4.1]{WilsonLines}}]\label{prop:unimodular}
$|K_\omega(i\tau)|=1$ for every real $\tau$ with $\xi(\frac12+\omega+i\tau)\neq0$.
\end{proposition}
\begin{proof}
$\xi(u)=\xi(1-u)$ gives $\xi(\frac12-\omega+i\tau)=\xi(\frac12+\omega-i\tau)$,
and $\xi(\bar u)=\overline{\xi(u)}$ gives
$\xi(\frac12+\omega-i\tau)=\overline{\xi(\frac12+\omega+i\tau)}$; numerator and
denominator of \eqref{eq:symbol} are conjugate.
\end{proof}

\begin{lemma}[Compression is multiplicative on causal kernels {\cite[Lemma~5.1]{WilsonLines}}]\label{lem:multiplicative}
If $\mathcal A,\mathcal B$ are causal convolutions and $f$ is supported in
$I_L$, then $P_L\mathcal A\mathcal BP_Lf=(P_L\mathcal AP_L)(P_L\mathcal BP_L)f$.
Consequently $H\mapsto P_L(h*\cdot)P_L$ is a homomorphism from the causal
convolution algebra onto a commutative algebra $\mathcal A_L$ of operators on
$L^2(I_L)$, in which all $V_{\omega,L}$ and $G_{\omega,L}$ lie.
\end{lemma}
\begin{proof}
$(I-P_L)\mathcal Bf$ is supported to the right of $I_L$; a causal operator
moves support only to the right; so $P_L\mathcal A(I-P_L)\mathcal BP_Lf=0$.
\end{proof}

\begin{proposition}[Flux and dichotomy {\cite[Propositions~4.5, 7.1, 7.2]{WilsonLines}}]\label{prop:dichotomy}
For fixed $\omega$, $L\mapsto\norm{V_{\omega,L}}$ is nondecreasing with limit
$\norm{V_\omega}_{L^2(\R)}$. If no zero of $\xi$ has $|\re\rho-\frac12|>\omega$,
then $K_\omega$ is inner on the right half-plane, $V_\omega$ is unitary and
causal on $L^2(\R)$, every $V_{\omega,L}$ is a contraction, and
\begin{equation}\label{eq:flux}
 \ip f{D_{\omega,L}f}=\norm{(I-P_L)V_\omega f}^2=\int_{L/2}^\infty|(V_\omega f)(x)|^2\dd x .
\end{equation}
If some zero has $|\re\rho-\frac12|>\omega$, then $\sup_L\norm{V_{\omega,L}}=\infty$.
\end{proposition}
\begin{proof}[Sketch]
Nested compressions; strong convergence $P_L\to I$. If no pole of $K_\omega$
lies in $\re p>0$ the realization contour may be moved to the imaginary axis,
where $K_\omega$ is unimodular, so $V_\omega$ is unitary and
$\norm f^2-\norm{P_LV_\omega f}^2=\norm{(I-P_L)V_\omega f}^2$, with the escaping
mass to the right by causality. A pole at $\re p=\delta-\omega>0$ contributes a
growing exponential to $k_\omega$, making $V_\omega$ unbounded, hence the
compressions unbounded in $L$.
\end{proof}

\subsection{Complete monotonicity, Bernstein functions, Blaschke factors}

A function $H$ on an interval $(p_0,\infty)$ is \emph{completely monotone} if
$(-1)^kH^{(k)}\geq0$ for all $k\geq0$. Bernstein's theorem states that $H$ is
completely monotone on $(p_0,\infty)$ if and only if
$H(p)=\int_{[0,\infty)}e^{-pt}\dd\mu(t)$ for a positive measure $\mu$ with the
integral finite for $p>p_0$; the measure is the causal inverse Laplace kernel
of $H$. Products of completely monotone functions are completely monotone, and
if $\Phi$ has a completely monotone derivative on $(p_0,\infty)$ then
$e^{-\Phi}$ is completely monotone there (after adding a constant, $\Phi$ is a
Bernstein function, and $e^{-\Phi}$ is then completely monotone by the standard
composition theorem; see \cite{SSV}). Both facts are used in
Section~\ref{sec:decomposition}.

A \emph{Blaschke factor} of the right half-plane is
$B_c(p)=\frac{p-c}{p+\bar c}$ with $\re c>0$; it is holomorphic and bounded on
$\re p>0$, unimodular on the axis, and its causal kernel is
$\delta_0-2\re c\,e^{-\bar ct}\mathbf 1_{t>0}$. A finite product of Blaschke
factors is the general rational function unimodular on the axis and equal to
$1$ at infinity. The \emph{exponential moving average} at rate $b>0$ is
$\EMA_b[g](x)=b\int_0^xe^{-b(x-y)}g(y)\dd y$, the causal convolution with
$be^{-bt}\mathbf 1_{t>0}$, of symbol $\frac b{p+b}$; so
$B_b=1-\frac{2b}{p+b}$ acts on kernels as $g\mapsto g-2\EMA_b[g]$.
